\section{Delzant's Theorem}

In this section we cover the following:
\begin{enumerate}
    \item What is the statement and proof of Delzant's Theorem
    \item Delzant in the $[-1, 1] = \Delta$ case
    \item Delzant's Theorem and uniqueness modulo toric equivariant symplectomorphism
\end{enumerate}
\subsection{Toric Actions and Moment Maps}

Toric actions are the most elementary Lie group actions because they are abelian. Let's see the definition of a toric action on complex $\mathbb{C}^{d}$-space:   
\begin{example}[Toric Action $T^{n} \circlearrowright \mathbb{C}^{n}$]
\label{ToricAction}
One common Lie group action is that of the torus, $T^{n}$ acting on $\mathbb{C}^{n}$: 
\begin{equation}
(e^{i \beta_{1}}, \dots, e^{i \beta_{n}}) \cdot (e^{i \theta_{1}}, \dots, e^{i \theta_{n}}) = (e^{i (\theta_{1} + \beta_{1})}, \dots, e^{i (\theta_{n} + \beta_{n})})
\end{equation}
The toric action is given by phase rotation in each of the complex coordinates. 
\end{example}

\noindent A toric action will also have a corresponding vector field. This is important to study as it will be used to find the moment map $\mu_{c}$

\begin{lemma}[Fundamental Vector Field $X^{\#}$ of Toric Action] 
Given a toric action as in \ref{ToricAction}, $$X^{\#}_{\theta_{1}, \dots \theta_{n}} = \frac{\partial}{\partial \theta_1} + \dots + \frac{\partial}{\partial \theta_n}$$
\end{lemma}

\begin{proof}
We take the derivative and observe:

\begin{align}
\left[ d \sigma_{\theta_{1}, \dots, \theta_{n}} \right]_{e} = (i e^{i \theta_{1}}, \dots, i e^{i \theta_{n}})
\end{align}

\noindent But this is the vector field tangent to $e^{i \theta_{1}}, e^{i \theta_{2}}, \dots e^{i \theta_{n}}$. This is, in local coordinates:

$$X^{\#}_{\theta_{1}, \dots \theta_{n}} = \frac{\partial}{\partial \theta_1} + \dots + \frac{\partial}{\partial \theta_n}$$
\end{proof}

Now, we must compute the symplectic form in polar coordinates. This is:

\begin{lemma}[A Symplectic Form on $\mathbb{C}^{n}$]
We choose a symplectic form:
$$\omega = \sum_{i} dz_{i} \wedge d\overline{z_{i}} = \sum_{i} |z_i| d \theta_{i} \wedge d|z_{i}|$$
\end{lemma}

\begin{proof}
In \ref{ComplexForm}, we proved that: $\omega_{\mathbb{C}^{k}} = \frac{-i}{2} \sum_{k} dz \wedge \overline{dz}$ was a symplectic form for $(\mathbb{C}^{k}, \omega_{\mathbb{C}^{k}})$. We choose a multiple of this form to be our form:

$$\omega = \sum_{j \geq 1} dz_{j} \wedge d\overline{z_{j}} $$

\noindent Then, we use polar coordinates $dz_{j} = |z_{j}| d \theta_{j}$, and find

$$\omega = \sum_{j \geq 1} |z_j| d \theta_{i} \wedge d|z_{j}|$$
\end{proof}

Now, with a symplectic form $\omega$, and a fundamental vector field, we compute the moment map $\mu_{c}$ on $\mathbb{C}^{d}$. This is:

\begin{lemma}[Moment Map for $(\mathbb{C}^{d}, \omega, T^{d})$]
\label{ComplexMomentMap}
\begin{equation}
\left (\mu (z_{1}, \dots, z_{n}) \right)_{j} = \frac{|z_{j}|^2}{2} + \lambda_{j}
\end{equation}

\noindent Where $\lambda_{j} \in \mathbb{R}^{*}$
\end{lemma}

\begin{proof}
According to $d \mu_{c} = \omega(X^{\#}, \cdot)$, we observe that: 
$\left (d \mu_{c} \right)_{j} = \left[ \omega(X^{\#}_{\theta_{1}, \dots \theta_{n}}, \cdot) \right]_{j} = |z_{j}| \cdot d|z_{j}|$. We integrate both sides to find: $\left (\mu \right )_{j} = \frac{|z_{j}|^2}{2} + \lambda_{j} \in (\mathbb{R}^{d})^{*}$
\end{proof}

\noindent \textbf{References:} \textcite{Guillemin1994} \textcite{Terek2020}, \textcite{Silva2001}

\subsection{Definitions Concerning the Delzant Construction}

\begin{remark}
Now, with basic results behind us for the toric action $T^{d} \circlearrowright M$, we can study the geometry of the image of the moment map as in \textcite{Delzant1988}. 
\end{remark}

Delzant's construction is an association between a symplectic manifold with toric action: $(M_{\Delta}, \omega_{\Delta}, T^{d})$ to each Delzant space $\Delta$, where $\Delta$ is a type of polytope, called a Delzant polytope. Delzant's construction links $M_{\Delta} \leftrightarrow \Delta$ via the Delzant moment map. Now, we will study the construction:

\begin{remark}
In this case, we denote the Delzant construction moment map: $J: M_{\Delta} \to (\mathbb{R}^{d})^{*}$, where as the moment map $\mu_{c}$ denotes the one here $\ref{ComplexMomentMap}$.
\end{remark}

\begin{defn}[Delzant Polytope]
A Delzant polytope $(\Delta, n, d)$ is a convex polytope in $(\mathbb{R}^{n})^{*}$. There are $d$-faces of a Delzant polytope $\Delta$, and $n$ edges meeting in each vertex $p$. The edges meeting at each $p$ are rational. Furthermore, $v_{1}, \dots, v_{n}$ can be chosen to be a basis of $(\mathbb{Z}^{n})^{*}$. Each one of these faces is defined by a normal equation, where $u_{i}$ is the inward pointing normal vector to each face \textcite{Delzant1988}

\begin{equation}
\langle u_{i}, x \rangle \geq \lambda_{i}, \; i \in \{1, \dots, d \}
\end{equation}
\end{defn}

\noindent Now, we define several maps used in the Delzant construction:

\begin{lemma}[Delzant Quotient Map]
Let $e_{1}, \dots, e_{d}$ be standard basis vectors of $\mathbb{R}^{d}$. There exists a linear map $\pi'$: 

$$\pi': T^{d} \to T^{n}$$
\end{lemma}

\begin{proof}
Let $e_{1}, \dots, e_{d}$ be standard basis vectors of $\mathbb{R}^{d}$ and consider the map:
\begin{equation}
\pi: \mathbb{Z}^{d} \to \mathbb{Z}^{n} 
\end{equation}
Consider its natural extension by inclusion:
\begin{equation}
\pi: \mathbb{R}^{d} \to \mathbb{R}^{n} 
\end{equation}
Then, there exists a quotient map: $\pi': T^{d} \to T^{n}$, such that the diagram commutes. The quotient map: $\pi'$ does the following:
\begin{equation}
\pi': [r] \mapsto [\pi(r)] 
\end{equation}
\begin{center}
\begin{tikzcd}
    \mathbb{Z}^{d} \arrow[r, "\pi"] \arrow[d, hook] & \mathbb{Z}^{n} \arrow[d, hook] \\
    \mathbb{R}^{d} \arrow[r, "\pi"] \arrow[d, two heads] & \mathbb{R}^{n} \arrow[d, two heads] \\
    T^{d} \arrow[r, "\pi'"] & T^{n}
\end{tikzcd}
\end{center}
\end{proof}

\begin{remark}
\label{NSpaceRemark}
It is an instructive exercise that $\pi'$ has a less obvious form. Let the coordinates of each normal vector be $u_{i} = (u_{i,1}, \dots, u_{i,n})$. Then,
\begin{equation}
\pi'(z_{1}, \dots, z_{d}) \mapsto (\prod_{i=1}^{d} z_{i}^{u_{i,1}}, \dots, \prod_{i=1}^{d} z_{i}^{u_{i,n}})
\end{equation}
\end{remark}

Since $\pi'$ is a linear map, we study its kernel, which we define here:

\begin{defn}[Delzant $N$ Space]
Let the kernel of $\pi'$ be $N$. 
\end{defn}

\begin{lemma}[Exact Sequence Defining $N$]
$\mathfrak{n}$ can equivalently be defined by the exact sequence:

$$0 \to \mathfrak{n} \to \mathbb{R}^{d} \to \mathbb{R}^{n} \to 0  $$
\end{lemma}
\begin{proof}
This follows from the exact sequence:
$$0 \to N \to T^{d} \to T^{n} \to 0$$
\noindent Now we take the differential of each map and move to the Lie algebra level:

$$0 \to \mathfrak{n} \to \mathbb{R}^{d} \to \mathbb{R}^{n} \to 0  $$
\noindent The conclusion follows
\end{proof}

\noindent \textbf{References}: \textcite{Silva2001} \textcite{Guillemin1994}, \textcite{Terek2020}

\subsection{Construction of Delzant Space $M_{\Delta}$ given $\Delta$}

Now, for a polytope $\Delta \in (\mathbb{R}^{n})^{*}$, we will construct a symplectic manifold $M_{\Delta}$, such that its moment map $J: M_{\Delta} \to \Delta$. First we need a few definitions:

\begin{defn}[Delzant Inclusion Map]
$$i: N \to T^{d}$$
\end{defn}

\begin{defn}[Dual Inclusion Map]
Define: 
\begin{align}
i^{*}: (\mathbb{R}^{d})^* \to \mathfrak{n}^{*} \text{ where } i^{*} \text{ is the derivative transpose of } i
\end{align}
\end{defn}

\begin{lemma}[Almost $M_{\Delta}$] 
The set $(i^{*} \circ J)^{-1}(0)$ is a compact subset of $\mathbb{C}^{d}$, \\
and $N$ acts freely on this set.
\end{lemma}

\begin{proof}
From the exact sequence:
\begin{equation}
0 \to \mathfrak{n} \to \mathbb{R}^{d} \to \mathbb{R}^{n} \to 0  
\end{equation}
we conclude that $\ker i^{*} = \operatorname{Im} \pi^{*}$. Furthermore, it is an observation that:  
\begin{equation}
\label{ImageofMDelta}
(i^{*} \circ \mu_{c})^{-1}(0) = \mu_{c}^{-1}(\pi^{*}(\Delta))
\end{equation}
Note that the stabilizer group of $z$ in $T^{d}$ is $T^{I}$, where $I$ is the number of $0$'s in the coordinate expansion of $z$. Note when $p$ is in the interior of $\Delta$, the number of $0$'s for a given $z_{i}$ is zero. This is because:
\begin{equation}
p \in \operatorname{Int} \Delta \Longleftrightarrow \langle u_{i}, p \rangle > \lambda_{i}
\end{equation}
When $p$ a vertex of the polygon, assume $p=0$, so that:
\begin{equation}
\langle u_{i}, p \rangle = 0, \; i \in \{1, \dots, n \}
\end{equation}
Then, we note that $\pi(T^{I}) \cong T^{n}$ and therefore, we have that $T^{I} \cap N = \emptyset$. Since $T^{I}$ is also the stabilizer group, we have that $N$ acts freely at $z$. Furthermore, note that since $J$ is proper, $(i^{*} \circ J)^{-1}(0)$ is compact.  
\end{proof}

\begin{remark}
In order to define the Delzant Space $M_{\Delta}$, we need to remove the free action of $N$ on $(i^{*} \circ \mu_{c})^{-1}(0)$. For this, we cite the following Theorem. 
\end{remark}

\begin{theorem}[Marsden-Weinstein] 
\label{Marsden-Weinstein}
Let $(M, \omega, G, \mu)$ be a Hamiltonian $G$-space and $p \in \mathfrak{g}^*$ be a regular value for $\mu$. Assume that the action $G_p \circlearrowright \mu^{-1}(p)$ is free and proper. Then the quotient $\mu^{-1}(p)/G_p$ has a unique symplectic form $\omega_0$ characterized by the relation $\pi^*\omega_0 = i^*\omega$, where $\pi: \mu^{-1}(p) \rightarrow \mu^{-1}/G_p$ is the quotient projection and $i: \mu^{-1}(p) \hookrightarrow M$ is the inclusion. We call $(\mu^{-1}(p)/G_p, \omega_0)$ the reduction of $(M, \omega)$ to level $p$. For more details, see \textcite{MarsdenWeinstein1974} and \textcite{Terek2020}
\end{theorem}

The following lemma is important as we reduce:

\begin{lemma}[Moment Map of Lie Subalgebra]
\label{MomentMapLieSubAlgebra}
Let $\mathfrak{g}$ be a Lie algebra and $\mathfrak{h} \subseteq \mathfrak{g}$ a Lie subalgebra. Suppose a symplectic manifold $(M, \omega)$ is equipped with a Hamiltonian action of a Lie group $G$ with moment map 
\[
\mu_c: M \to \mathfrak{g}^*.
\]
Let $H \subseteq G$ be a Lie subgroup with Lie algebra $\mathfrak{h}$, and consider the restricted action $H \circlearrowright M$. Then the dual of the inclusion map $i: \mathfrak{h} \hookrightarrow \mathfrak{g}$ induces a projection
\[
i^*: \mathfrak{g}^* \to \mathfrak{h}^*,
\]
and the composition
\[
\mu_H := i^* \circ \mu: M \to \mathfrak{h}^*
\]
is a moment map for the $H$-action.
\end{lemma}

\begin{proof}
Since Hamiltonian functions arising from the restricted action $H \circlearrowright M$ are restrictions of Hamiltonian functions arising from $G \circlearrowright M$, it follows that:
\begin{align}
\mu_{H, c}: \mathfrak{h} \to C^{\infty}(M) &&  \mu_{c}: \mathfrak{g} \to C^{\infty}(M) 
\end{align}
\noindent are related by $\mu_{H, c} = \mu_{c} \circ i$. Then, we note that: $\langle \mu, i \circ X \rangle = \mu_{H,c}(X)$. Therefore, $\langle i^{*} \circ \mu, X \rangle = \mu_{H,c}(X)$, and we note that: $\mu_{H} = i^{*} \circ \mu$. 
\end{proof}

\begin{theorem}[Construction of $M_{\Delta}$]
Given $\Delta \in( \mathbb{R}^{n})^{*}$ a Delzant Polytope, there exists: 
$(M_{\Delta}, \omega_{0}, T^{n})$ such that $J(M_{\Delta}) = \Delta$
\end{theorem}

\begin{proof}
From the Marsden-Weinstein Theorem \ref{Marsden-Weinstein}, 
we reduce $\mathbb{C}^{d}$ with respect to $N$. Since $N$ acts freely on $(i^{*} \circ J)^{-1}(0)$, the reduced space is a smooth symplectic manifold, which we declare $M_{\Delta}$. The symplectic manifold $M_{\Delta}$ admits a $T^{n}$ action after reduction. Define $\mu_{c}$ to be the moment map $\mu_{c}$ arising from the action of $T^{n}$ on $M_{\Delta}$. Consider the inclusion $i': \mathbb{R}^{n} \to \mathbb{R}^{d}$ and its transpose, $(i')^{*}: (\mathbb{R}^{d})^{*} \to (\mathbb{R}^{n})^{*}$. By \ref{MomentMapLieSubAlgebra}, the moment image of $M_{\Delta}$ is by $\ref{ImageofMDelta}$, $J(M_{\Delta}) = \Delta$
\begin{align}
J = (i')^{*} \circ \mu_{c} && M_{\Delta} = (i^{*} \circ \mu_{c})^{-1}(0)/N && J(M_{\Delta}) &\cong \Delta \in \mathbb{R}^{n}^{*}
\end{align} 
\end{proof}

\noindent \textbf{References}: \textcite{Silva2001} \textcite{Guillemin1994}, \textcite{Terek2020}

\subsection{Delzant Construction for $\Delta = [-1,1]$}

We will do the Delzant Construction for $\Delta = [-1,1]$. 

\begin{question}[What is $N$ given $\Delta$]
Claim: $N = \{(t, t) : t \in [0, 2\pi)\}$.
\end{question}
\begin{proof}
When $\Delta = [-1,1]$, the normal vectors are $u_{1} = -1$ and $u_{2} = 1$. In this case, $N = \ker \pi'$ where:
\begin{equation}
\pi'(z_{1}, z_{2}) = z_{1}^{-1} z_{2}
\end{equation}
\noindent We use Remark \ref{NSpaceRemark}, and observe that $N = \{(t, t) : t \in [0, 2\pi)\}$.
\end{proof}
\begin{question}[What is $\mu$ given $\Delta$]
Claim: $J(z) = \frac{(|z_{1}|^2, |z_{2}|^2)}{2} + (-1, -1) $
\end{question}
\begin{proof}
Since the interval $\Delta$ can also be seen as $x \geq -1$ and $-x \geq -1$ we have that $\lambda_{1} = -1$ and $\lambda_{2} = -1$.
Using \ref{ComplexForm}, the complex space $\mathbb{C}^{2}$ has moment map: $\mu(z) = \frac{(|z_{1}|^2, |z_{2}|^2)}{2} + (-1, -1)  $
\end{proof}
\begin{question}[What is $J$ given $\Delta$]
Claim: $J = i^{*} \circ J = \frac{|z_{1}|^2 + |z_{2}|^2}{2} - 2$
\end{question}
\begin{proof}
Note that $i: \mathfrak{n} \to \mathbb{R}^2$, given by $v \mapsto (v,v)$. Note that $\langle i^{*}(z_{1}, z_{2}), v \rangle = \langle (z_{1}, z_{2}), (v, v) \rangle = (z_{1} + z_{2}) v$. We use the property that $\langle a^{*} w, v \rangle = \langle w, av \rangle$.
Therefore, $i^{*}(z_{1}, z_{2}) = z_{1} + z_{2}$, and we note that:
$$i^{*} \circ J = \frac{|z_{1}|^2 + |z_{2}|^2}{2} - 2$$
\end{proof}
\begin{question}[What is $M_{\Delta}$]
Claim: $M_{\Delta} \cong \mathbb{C} P^{1}$
\end{question}
\begin{proof}
We note that $N \cong U(1)$. Therefore, $X_{\Delta} \cong S^{3}/N \cong S^{3}/U(1) \cong \mathbb{CP}^{1}$. 
\end{proof}

\subsection{Convexity of Moment Maps}
For the other direction, we use the \textcite{AtiyahGS1982}. Following the treatment in \textcite{Silva2001}, we sketch only one part of the proof. 
\begin{theorem}[Atiyah-Guillemin-Sternberg] 
Let \((M, \omega)\) be a compact connected symplectic manifold, and let \(T^{m}\) be an \(m\)-torus. Suppose that \(\psi: T^{m} \to \operatorname{Symp}(M, \omega)\) is a Hamiltonian action with moment map \(\mu: M \to \mathbb{R}^{m}\). Then:
\begin{enumerate}
    \item the levels of \(\mu\) are connected;
    \item the image of \(\mu\) is convex;
    \item the image of \(\mu\) is the convex hull of the images of the fixed points of the action.
\end{enumerate}
\end{theorem}
Consider the statements:
\begin{enumerate}
    \item $A_{n}$ - the levels of $\mu$ are connected for any $T^{n}$ action
    \item $B_{n}$ - the image of $\mu$ is convex for any $T^{n}$ action
\end{enumerate} From Morse Theory, we have that $A_{1}$ and $B_{1}$ hold.

\begin{theorem}[Part of Proof of AGS]
$A_{m-1}$ implies $B_{m}$
\end{theorem}

\begin{proof}
\noindent Consider a moment map $\mu$ from a toric action $T^{n} \circlearrowright M$. We can consider the toric action of an $(n-1)$-dimensional sub-torus on $M$, $T^{n-1} \circlearrowright M$, by using an injective matrix $A \in \mathbb{Z}^{n \times (n-1)}$ and defining:
\begin{align}
\psi_{A}: T^{n-1} &\to \operatorname{Symp}(M, \omega) && \theta \mapsto \psi_{A\theta}
\end{align}

\noindent Since $A^{T}: \mathbb{R}^{n} \to \mathbb{R}^{n-1}$ is a linear map from the Lie algebra of the torus onto the Lie algebra of its subgroup, the moment map of $\psi_{A}$ is $\mu_{A} = A^{T} \mu$. Consider $p, p_{0} \in \mu_{A}^{-1}(\xi)$, then $A^{T} \mu(p) = \xi = A^{T} \mu(p_{0})$. Now it follows that:

\begin{equation}
\mu_{A}^{-1}(\xi) = \{p \in M: \mu(p) - \mu(p_{0}) \in \operatorname{Ker}(A^{T}) \}
\end{equation}
\noindent Since $\mu_{A}^{-1}(\xi)$ is a connected manifold for regular value $\xi$ and is locally path connected, the inverse image is path connected. Therefore, for a path $t \mapsto p_{t} \in \mu_{A}^{-1}(\xi)$, there also exists a path $t \mapsto \mu(p_{t}) - \mu(p_{0}) \in \operatorname{Ker}(A^{T})$. Since $\operatorname{Ker}(A^{T})$ is $1$-dimensional, it follows that all elements $\mu(p_{t})$ must be a linear superposition of $\mu(p_{0})$ and $\mu(p_{1})$ - two points determine a line!

\begin{equation}
(1-t) \mu(p_{0}) + t \mu(p_{1}) \in \mu(M), \; t \in [0,1]
\end{equation}
Choosing points $p_{0}$ and $p_{1}$ closely approximated by rational $p_{0}'$ and $p_{1}'$ we run this argument again - finding a matrix which contains $\mu(p_{0}') - \mu(p_{1}')$ in its kernel. Taking limits, we find that $\mu(M)$ is convex.
\end{proof}

\noindent \textbf{References}: \textcite{Silva2001}, \textcite{Delzant1988}, \textcite{Terek2020}

\subsection{Uniqueness of Delzant Spaces}
Now, we have shown that for each Delzant polytope, there exists a corresponding Delzant space, and given a Delzant space, we may find the polytope. Actually, the association of $\Delta$ and $M_{\Delta}$ is unique, up to toric-equivariant symplectomorphism as in  \textcite{Delzant1988}.

\begin{defn}
Given a Delzant polytope $\Delta$, if there exist moment maps $\mu_1$ and $\mu_2$, and Delzant spaces, $X_{\Delta, 1}$ and $X_{\Delta, 2}$, such that there exists an isomorphism $\gamma: \operatorname{Im}(\mu_1) \to \operatorname{Im}(\mu_2)$ and $\phi: X_{\Delta, 1} \to X_{\Delta, 2}$ is a $\gamma$-equivariant symplectomorphism, then $\phi$ is called a $\gamma$-equivariant symplectomorphism, which is a diffeomorphism in the symplectic category that respects the moment maps.
\end{defn}
\begin{center}
\begin{tikzcd}
X_{\Delta, 1} \arrow[r, "\phi"] \arrow[d, "\mu_1"'] & X_{\Delta, 2} \arrow[d, "\mu_2"] \\
\operatorname{Im}(\mu_1) \arrow[r, "\gamma"'] & \operatorname{Im}(\mu_2)
\end{tikzcd}
\end{center}

\noindent Now, we would like to extend these results to the non-Abelian case $G \circlearrowright M$. It turns out that so long as the action of $G$ has regular values, the uniqueness of the moment map image is determined by the toric (Abelian) part of $G$, which is known as the maximal torus of $G$. 

\noindent \textbf{References: } \textcite{Blaom1996}, \textcite{GuilleminSternberg1982}.

\newpage

\setcounter{section}{4}